\subsection{What this experiment excludes, and what it does not}
\label{sec:excludes}

\noindent\emph{Excluded.} A continuously transmitting, unresolved carrier
lying outside the molecular-line mask, inside the searched frequency range
and inside the linear-drift grid, and present during one of the archival
epochs, would have been found nine times in ten above the
$\mathrm{EIRP}_{90}$ of \S\ref{sec:pninetysel} --- a median of
$\EirpNinetySysMedA$\,W per system on its best window. No such carrier is
present toward any of the \NSystems{} systems searched.

\noindent\emph{Not excluded.} A transmitter absent during the archival
observation, or on for less than about a tenth of the time: at a duty cycle
of one tenth, recovery falls to \DwellFineT{} per cent for the fine class
and \DwellCoarseT{} per cent for the coarse
(Appendix~\ref{app:dwelldetail}), and those trials are undrifted, so the
joint drifting-and-intermittent case is not measured at all. Emission inside
the $\pm\MaskHalfKms$\,km\,s$^{-1}$ molecular-line mask, at any strength.
Drift outside the searched grid, or drift curved enough to leave a channel
within the track. Signals broad enough to be removed with the continuum
baseline. Anything below a window's own completeness. And any emission from
the \NGaiaCensus{} stars within 40\,pc that ALMA has not observed, which is
almost all of them.

\subsubsection{The survey's on-source time, in both directions}
\label{sec:exposure}

Two bookkeeping biases of opposite sign affect the exposure, and neither may
be quoted without the other. The extractor retained a single \texttt{FIELD}
row per execution block, so in \TruncNBlocks{} of the \ExpNBlocks{} released
blocks --- \TruncNWin{} windows toward \TruncNStars{} stars --- further
in-beam pointings on the same star were discarded, \TruncHoursLost\,h of
real on-source time, \TruncPctSurvey{} per cent of the survey. Nothing
published is thereby wrong, because every noise level was measured from the
data actually used; but the \emph{achievable} sensitivity in those windows
is understated by a median factor \TruncGainMed{} in flux density, and those
blocks were re-extracted and re-searched. In the opposite direction, the
released on-source time is the number of integrations times the median dump
length rather than a sum of dump lengths, so short sub-scan boundary dumps
are over-counted; the exact computation moves the survey total from
\OcOldMax{} to \OcNewMax\,h. \emph{The first bias costs the survey time it
had and the second credits it with time it never had}, and they offset to
within \ExpNetLo--\ExpNetHi\,h.

\subsubsection{The omitted annual parallax} \label{sec:parallax}

The extraction phase-rotates the visibilities to each star's
\emph{barycentric} direction while the visibilities see the apparent one, so
annual parallax is omitted from the assumed phase centre. The resulting
amplitude loss is the normalised dirty beam evaluated at the parallax
displacement in units of the synthesised beam: parameter-free,
multiplicative and shape-preserving. Over the \PxNWin{} windows carrying the
astrometric keys it is \PxLossMedPct{} per cent in the median window and
\PxLossMaxPct{} per cent at worst, with \PxNGtTen{} windows losing more than
ten per cent. It is negligible throughout the ACA data that dominate this
survey and matters only for nearby stars observed in long-baseline 12\,m
configurations. \emph{We state it and do not apply it.} The sign is
optimistic --- every limit quoted here is too deep by $1/(1-{\rm loss})$ in
its own window --- and the completeness campaign cannot see the term,
because its carriers are deposited at the same assumed position the
estimator evaluates. Correcting it would leave the number of systems
reaching $10^{15}$\,W unchanged at \PxNSysCorr, because a long-baseline
block is never a system's most sensitive window.

\subsubsection{Two blind spots}

The search is Stokes~$I$ only. A fully polarised carrier is not missed, but
the discrimination that polarisation would offer against unpolarised
astrophysical emission is unavailable. ALMA's feeds are linear, so a
per-hand search would help only a carrier aligned with one feed; a
circularly polarised signal divides equally between XX and YY, and its
circularity is in the cross-hands, which this analysis does not extract.
The frequency coverage, finally, is the archive's and not a choice. None of
the ``magic frequency'' arguments that have guided centimetre-band searches
--- the 1420\,MHz line, its multiples, the water hole, or $\pi$ and $e$
times them --- falls in an ALMA band, so a transmitter deliberately placed
on one lies outside this experiment by construction and its absence here is
not evidence against it.

%% ---------------------------------------------------------------------
%% REQUEST (results owner -> numbers owner).  This file uses
%% \EirpNinetySysMedA, requested in 05_results.tex.
%%
%% REQUEST (results owner -> integrator).  Changes made here that touch
%% other files:
%%   - The sensitivity material (P90, the injection campaign, the budget
%%     discussion) has moved to section 5.1 in 05_results.tex, which now
%%     carries \label{sec:pninetysel}.  This file keeps only
%%     \label{sec:excludes}.  The two labels were previously on one
%%     subsection, which made every \ref to one of them resolve to the
%%     other.
%%   - fig:dwellsel is deleted here; the float belongs with the injection
%%     material in Appendix app:inject / app:dwelldetail, which is where
%%     the numbers now point.  Nothing outside this file cited it.
%%   - sec:exposure and sec:parallax are now numbered subsubsections
%%     rather than starred paragraphs, so \ref to them prints their own
%%     number instead of the enclosing subsection's.
%% ---------------------------------------------------------------------
